\documentclass[11pt,a4paper]{article}
\usepackage{amsmath,amssymb,amsthm}
\usepackage[margin=2.5cm]{geometry}
\usepackage{hyperref}

\newtheorem{definition}{Definition}[section]
\newtheorem{theorem}[definition]{Theorem}
\newtheorem{proposition}[definition]{Proposition}
\newtheorem{lemma}[definition]{Lemma}
\newtheorem{corollary}[definition]{Corollary}
\newtheorem{conjecture}[definition]{Conjecture}
\newtheorem{remark}[definition]{Remark}

\title{Hyperseed Formalization:\\
  Hyperseed v2}
\author{ProtomegaTron (automated formalization)}
\date{2026-09-18}

\begin{document}
\maketitle

\begin{abstract}
We formalize the intellectual structure of Ben Goertzel's Substack article
``Hyperseed v2'' (2026-03-05) within the Hyperseed ontology itself. This
article IS the Hyperseed ontology --- it describes the fiber bundle framework,
semantic primitives, d-calculus, and adaptive base space that all other
formalizations use. The self-referential formalization maps the framework's
own concepts to the framework's own structures: concepts as fibers over a
primitive base, d-calculus as fiber transport, curvature as semantic
obstruction, and inference acceleration as fiber-base alignment efficiency.
\end{abstract}

\tableofcontents

%% =================================================================
\section{The base space: semantic primitives}
\label{sec:base}
%% =================================================================

\begin{theorem}[Semantic primitives as base space --- claims 1--5, 27]
\label{thm:base}
The quest for semantic primitives traces from Leibniz's Universal
Characteristic (factoring concepts into prime elements) through Carnap's
formal reduction (logical combinations of universals), Wierzbicka's Natural
Semantic Metalanguage (concepts occurring as single words in every human
language), and Chalmers' Constructing the World (core concepts + raw
observations + smart reasoner = understanding of any concept).

Hyperseed v2 fills the practical AI gap: none of these philosophical
programs has been worked out at the level needed to drive AI inference.

In the Hyperseed fiber bundle, the base space $B$ is the set of semantic
primitives:
\[
  B = \{p_1, p_2, \ldots, p_n\} \qquad \text{(semantic primitives)}
\]

All other concepts are fibers over this base. A concept $c$ is a
fiber structure --- a pattern of relationships among primitives:
\[
  c \in F_b \qquad \text{for some } b \in B^k
  \text{ (combination of primitives)}
\]

Reducing a concept to primitives is a fiber-bundle projection $\pi$:
\[
  \pi: E \to B \qquad c \mapsto (p_{i_1}, \ldots, p_{i_k})
\]

The projection preserves structural relationships: the fiber structure
of a complex concept encodes not just which primitives are involved but
how they relate to each other. This is compositional semantics: the
meaning of a complex concept is determined by the fiber structure of its
components and their composition rules.

The specific primitive set includes both conventional and eccentric
choices. The exact list matters less than the compositional structure ---
it is not a reference taxonomy or philosophical curiosity but a tool
designed to change what a reasoning system does.
\end{theorem}

%% =================================================================
\section{Fiber transport: the d-calculus}
\label{sec:dcalculus}
%% =================================================================

\begin{definition}[d-calculus as fiber transport --- claims 11--14, 28]
\label{def:dcalculus}
The d-calculus formalizes relationships between fibers --- how concepts
relate to each other through their fiber structures. It provides the
mathematical language for computing with fiber bundles.

\textbf{Connection:} A connection $\nabla$ provides maps that transport
meaning along paths in the base space. This is how inference traverses
concept space: by following fiber-preserving paths:
\[
  \nabla_\gamma: F_{b_1} \to F_{b_2} \qquad
  \text{(parallel transport along path } \gamma: b_1 \to b_2\text{)}
\]

When we say ``concept X in article Y maps to fiber structure Z,'' we're
using d-calculus transport operations to move from one fiber (the
article's conceptual space) to another (the Hyperseed primitive space).

\textbf{Curvature:} When transporting a concept along different paths
in the base space yields different results, there's curvature --- a
semantic obstruction:
\[
  R(\gamma_1, \gamma_2) = \nabla_{\gamma_1} - \nabla_{\gamma_2}
  \neq 0 \implies \text{genuine conceptual complexity}
\]

Curvature indicates that the concept cannot be reduced away --- there's
irreducible structure that depends on path (context). This is the
mathematical formalization of ``context matters for meaning.''

\textbf{Holonomy:} Going around a closed loop in concept space and
arriving at a different meaning than you started with:
\[
  \text{Hol}(\gamma) = \nabla_\gamma \circ \nabla_{\gamma^{-1}} \neq
  \text{id} \implies \text{concept drift}
\]

Holonomy captures concept drift and contextual meaning shift: the
same word, after being used in a sequence of contexts, may have shifted
meaning. The holonomy measures how much it shifted.
\end{definition}

%% =================================================================
\section{Adaptive base: evolving primitives}
\label{sec:adaptive}
%% =================================================================

\begin{proposition}[Base space evolution --- claims 15--18, 29]
\label{prop:adaptive}
The collection of primitives is not fixed from the start --- it can be
adaptively sculpted over time as the system learns and encounters new
domains. The base space itself evolves:
\[
  B(t) = B(t-1) \cup \{p_{\text{new}}\} \setminus \{p_{\text{retired}}\}
\]

New primitives emerge when existing ones can't efficiently decompose new
concepts (the fiber structure over the current base is too complex),
and old ones are retired when they become redundant (no concepts use
them as essential base components).

The system meta-learns: it learns not just new concepts (new fibers)
but learns which primitives are most useful for learning new concepts
(optimizes the base). The primitive set is a learned representation.

There's a wide range of reasonable primitive sets --- tragically bad
decisions are possible but many good options exist. The adaptive
mechanism finds good ones through experience. The base space that works
for physics articles may need extension for ethics articles; the
base space for everyday reasoning may differ from the base space for
mathematical reasoning.

In Hyperseed terms, the base is not a static mathematical object but
a dynamic one that co-evolves with the fiber. The fiber content
(concepts learned) feeds back into base structure (which primitives
are available), which affects future fiber construction (how new
concepts are decomposed). This is a fiber-base feedback loop that
drives the system toward increasingly efficient decomposition.
\end{proposition}

%% =================================================================
\section{Inference acceleration: fiber-base alignment}
\label{sec:acceleration}
%% =================================================================

\begin{theorem}[Semantic highways --- claims 19--22, 30]
\label{thm:highways}
Hyperseed primitives provide ``semantic highways'' --- pre-computed
reduction paths that inference can follow rather than exploring concept
space from scratch:
\[
  \text{Highway}(c_1, c_2) = \text{path via base: }
  \pi(c_1) \to \pi(c_2) \text{ then lift}
\]

Instead of searching for a connection between two concepts in the full
fiber space (exponential), the system projects both to base (fast),
finds a path between their primitive decompositions (fast, because
the base is small), and lifts back to concept space (fast, because
the lift is fiber-preserving).

The speedup is exponential for cases requiring rigor, imagination, and
cross-domain thinking: the difference between exploring all paths in
concept space and following semantic highways through primitive space.

Cross-domain reasoning becomes natural: concepts from different domains
that reduce to similar primitive patterns are discovered to be related.
A physics concept and an ethics concept that share a primitive
decomposition pattern are structurally analogous --- and the system
discovers this automatically through the base-space projection.

The practical target is inference control in commonsense or scientific
reasoning contexts --- exactly the areas where LLMs struggle most.
LLMs lack the semantic primitive infrastructure that would enable
systematic cross-domain reasoning; they rely on surface-level pattern
matching rather than primitive-level structural analysis.

In Hyperseed terms, the inference speedup comes from fiber-base
alignment: when concept and primitive are well-aligned (the fiber
is ``thin'' --- close to the base projection), inference follows the
fiber-preserving path (fast). When misaligned (the fiber is ``thick''
--- far from any base projection), inference must search (slow).
Semantic highways are high-alignment paths through the bundle.
\end{theorem}

%% =================================================================
\section{Observer-relative ontology}
\label{sec:relative}
%% =================================================================

\begin{proposition}[No fixed ground truth --- claims 23--25]
\label{prop:relative}
The ontology is observer-relative: the primitive decomposition of a
concept depends on who is decomposing it. Different observers with
different primitive sets produce different but equally valid
decompositions:
\[
  \pi_1(c) \neq \pi_2(c) \qquad \text{both valid, observer-relative}
\]

There is no fixed ground truth about what the ``real'' primitives are.
The primitives are tools, not discoveries. Different tasks call for
different primitive sets. The philosophy has a postmodern, relativist,
Buddhist flavor expressed with analytical-philosophy-level mathematical
precision.

This connects to the observer-relative quantumity framework from
``Sometimes a Smaller System Should Model a Bigger System as Quantum'':
just as quantumity is observer-relative, semantic decomposition is
observer-relative. The base space itself is not absolute --- it depends
on the observer (the decomposing agent). Different agents with different
cognitive architectures will naturally develop different primitive sets.

The mathematical precision is important: ``observer-relative'' does not
mean ``anything goes.'' Each primitive set must satisfy structural
constraints (compositional consistency, completeness for the target
domain, computational tractability). Within those constraints, many
valid choices exist. The relativism is constrained, not arbitrary.
\end{proposition}

%% =================================================================
\section{Cross-references}
%% =================================================================

\begin{remark}[Self-reference and other connections]
This formalization is unique among all Hyperseed formalizations because
it is self-referential: the article being formalized IS the framework
used for formalization. The interpretation of ``Hyperseed v2'' in terms
of Hyperseed ontology is the identity map --- every concept in the
article is already a Hyperseed concept.

Connections to other formalizations:
\begin{itemize}
\item \textbf{Every other article}: uses the framework described here.
  The fiber bundle structure, d-calculus, semantic primitives, and
  adaptive base space are the machinery applied to all other articles.
\item \textbf{Evidence Is to Logic What Energy Is to Physics} (2026-03-10):
  evidence conservation. The quantale Noether theorem provides the
  evidence conservation principle that constrains fiber operations
  within the Hyperseed bundle.
\item \textbf{Sometimes a Smaller System Should Model a Bigger System
  as Quantum} (2026-03-18): observer-relative quantumity. Extends the
  observer-relative ontology to quantumity itself.
\item \textbf{Seeding RSI} (2026-07-28): architecture. Hyperseed v2's
  adaptive primitives connect to RSI: the system improving its own
  primitive set is a form of self-improvement.
\end{itemize}
\end{remark}

\end{document}
